
\subsection{命题~\ref{prop:7-1-1}、\ref{prop:7-1-2} 证明的收束}
\label{subsec:7-1-4}

前两个小节已经描述命题~\ref{prop:7-1-1}、
\ref{prop:7-1-2} 所述仓西结构构造的分析部分。
本小节将从整体上构造仓西结构，完成证明。
此证明基本上是对 \textbf{[FuOn99II]} 中论述的小幅修改。

先精确定义\liftIndex{I374}
$\mathcal M^{\mathrm{main}}_{k+1,m}(L;\beta)$
上的拓扑。
它类似于 \textbf{[FuOn99II]} 第~10 节所定义的拓扑，
具体如下。命题~\ref{prop:7-1-1}、
\ref{prop:7-1-2} 只涉及 $m=0$，
但研究球面气泡等现象时必须纳入 $m\ne0$。
本节只考虑 $k+1>0$；
$k+1=0$ 的情形将在第~7.4.1 小节讨论。

设
$(\Sigma_i,\vec z_i,\vec z_i^{\, +})
\in\mathcal M^{\mathrm{main}}_{k+1,m}$
是一列亏格零、带标记点的有边稳定曲线，
其中
$\vec z_i=(z_0,\ldots,z_k)$
为边界标记点，
$\vec z_i^{\, +}$ 为内部标记点；
参见定义~2.1.17 和定义~2.1.18。
又设 $w_i\colon\Sigma_i\to M$
是一列伪全纯映射，且
\[
\lim_{i\to\infty}
(\Sigma_i,\vec z_i,\vec z_i^{\, +})
=
(\Sigma_\infty,\vec z_\infty,\vec z_\infty^{\, +})
\in\mathcal M^{\mathrm{main}}_{k+1,m}.
\]
必要时取子列，可假设这些曲线的组合型不依赖 $i$。
对于 $\Sigma_\infty$ 中每个
不对应于 $\Sigma_i$ 中奇点的奇点，
取它的 $\epsilon$-邻域，
以 $V_{\epsilon,\infty}\subset\Sigma_\infty$
表示这些邻域之并。
在 $\Sigma_i$ 中取相应矩形区域之并
$V_{\epsilon,i}$。
于是有微分同胚
\[
u_{i,\epsilon}\colon
\Sigma_i\setminus V_{\epsilon,i}
\xrightarrow{\ \cong\ }
\Sigma_\infty\setminus V_{\epsilon,\infty}.
\]
参见 \textbf{[FuOn99II]} 第~7 节；
该节讨论的情形与这里类似，
而借助本书第~2.1.2 小节的论证
可将这里的情形化归过去。
见图~\ref{fig:7-1-7}\translatornote{原书图 7.1.7 的右侧邻域印作
$V_{\infty,i}$；据同页正文的定义，此处应为
$V_{\epsilon,\infty}$。重绘图采用正文记号，见疑误 A-0424。}。

\par\medskip\noindent\begin{minipage}{\textwidth}
\centering
\setlength{\unitlength}{0.74mm}
\begin{picture}(145,54)
\linethickness{0.55pt}
\put(16,35){\circle{18}}
\put(43,26){\circle{27}}
\put(81,25){\circle{29}}
\put(39,8){\circle{13}}
\put(82,5){\circle{12}}
\put(39,14){\line(0,1){2}}
\put(82,11){\line(0,1){1}}
\qbezier(24,31)(28,38)(31,35)
\qbezier(23,28)(27,23)(30,22)
\qbezier(56,29)(61,33)(67,29)
\qbezier(55,22)(61,19)(67,21)
\put(24,23){\circle{11}}
\put(61,25){\circle{12}}
\put(95,25){\circle*{1.5}}
\put(98,25){$\scriptstyle z_0$}
\put(17,48){$\scriptstyle V_{\epsilon,i}$}
\put(39,1){$\scriptstyle\Sigma_i$}
\put(104,39){\circle{12}}
\put(120,27){\circle{15}}
\put(139,22){\circle{13}}
\put(117,10){\circle{10}}
\put(139,6){\circle{10}}
\put(117,15){\line(0,1){5}}
\put(139,11){\line(0,1){5}}
\put(110,36){\line(1,-1){4}}
\put(127,26){\line(1,-1){5}}
\put(113,33){\circle*{1.5}}
\put(132,21){\circle*{1.5}}
\put(145,22){\circle*{1.5}}
\put(139,35){$\scriptstyle V_{\epsilon,\infty}$}
\put(145,23){$\scriptstyle z_0$}
\put(116,1){$\scriptstyle\Sigma_\infty$}
\end{picture}
\par\smallskip
{\renewcommand{\thefigure}{7.1.7}\refstepcounter{figure}
{\normalfont\bfseries 图~\thefigure}\label{fig:7-1-7}}
\end{minipage}\par\medskip

\begin{definition}{7.1.39}\label{def:7-1-39}
称
\[
\lim_{\mathrm{s},\,i\to\infty}
\bigl((\Sigma_i,\vec z_i,\vec z_i^{\, +}),w_i\bigr)
=
\bigl((\Sigma_\infty,\vec z_\infty,\vec z_\infty^{\, +}),
w_\infty\bigr),
\]
如果对每个 $\delta>0$，
存在 $\epsilon(\delta)>0$ 与 $I(\delta)$，
使得对 $i>I(\delta)$ 和
$0<\epsilon<\epsilon(\delta)$
以下条件成立：
\begin{enumerate}[label=\textup{(7.1.40.\arabic*)},leftmargin=6em]
\item\label{eq:7-1-40-1}
$\displaystyle
\lim_{i\to\infty}
(\Sigma_i,\vec z_i,\vec z_i^{\, +})
=(\Sigma_\infty,\vec z_\infty,\vec z_\infty^{\, +})
\in\mathcal M^{\mathrm{main}}_{k+1,m}$。
\item\label{eq:7-1-40-2}
$\displaystyle
\sup_{\Sigma_i\setminus V_{\epsilon,i}}
\operatorname{dist}(w_i,w_\infty\circ u_{i,\epsilon})
<\delta$。
这里先在 $M$ 上固定一个黎曼度量，
$\operatorname{dist}$ 是它诱导的距离。
\item\label{eq:7-1-40-3}
$V_{\epsilon,i}$、
$V_{\epsilon,\infty}$ 的每个连通分支
在 $w_i$、$w_\infty$ 下的像的直径，
分别都小于 $\delta$。
\end{enumerate}
\end{definition}

现在将定义~\ref{def:7-1-39}
推广到以下情形：亏格零带标记点的有边黎曼曲面
$(\Sigma_i,\vec z_i,\vec z_i^{\, +})$ 与
$(\Sigma_\infty,\vec z_\infty,\vec z_\infty^{\, +})$
本身可以不稳定，
但连同映射所得的
$((\Sigma_i,\vec z_i,\vec z_i^{\, +}),w_i)$ 与
$((\Sigma_\infty,\vec z_\infty,\vec z_\infty^{\, +}),w_\infty)$
是稳定的。
办法是添加标记点，使域稳定。
取保序单射
$\rho\colon\{1,\ldots,k'\}\to\{1,\ldots,k\}$
及单射
$\rho'\colon\{1,\ldots,m'\}\to\{1,\ldots,m\}$。
在 $k\ge0$ 时如下定义遗忘映射\liftIndex{I230}：
\[
\operatorname{Forget}_{\rho,\rho'}\colon
\mathcal M^{\mathrm{main}}_{k+1,m}(L;\beta)
\longrightarrow
\mathcal M^{\mathrm{main}}_{k'+1,m'}(L;\beta).
\]
给定
$((\Sigma,\vec z,\vec z^{\, +}),w)
\in\mathcal M^{\mathrm{main}}_{k+1,m}(L;\beta)$，
先考虑
\[
\bigl(
(\Sigma,
(z_0,z_{\rho(1)},\ldots,z_{\rho(k')}),
(z^+_{\rho'(1)},\ldots,z^+_{\rho'(m')})),w
\bigr).
\]
所得映射未必稳定。
若有不稳定分支 $D$，
它只可能属于以下六种情形，
而 $w$ 在该分支上必为常值映射；
参见图~\ref{fig:7-1-8}--\ref{fig:7-1-13}：
\begin{enumerate}[label=\textup{(7.1.41.\arabic*)},leftmargin=6em]
\item\label{eq:7-1-41-1} 球面，有一个奇点、没有标记点；
\item\label{eq:7-1-41-2} 球面，有两个奇点、没有标记点；
\item\label{eq:7-1-41-3} 球面，有一个奇点、一个标记点；
\item\label{eq:7-1-41-4} 圆盘，有一个边界奇点、没有标记点；
\item\label{eq:7-1-41-5} 圆盘，有两个边界奇点、没有标记点；
\item\label{eq:7-1-41-6} 圆盘，有一个边界奇点、一个边界标记点。
\end{enumerate}
注意 $k+1>0$，
因此任何圆盘分支 $D$ 都至少有一个标记点或奇点。
于是，如果它还有一个内部奇点，就必为稳定分支。

在 (7.1.41.1) 或 (7.1.41.4) 中，
删去 $D$；
在 (7.1.41.2) 或 (7.1.41.5) 中，
删去 $D$，再将原来与它相连的两个分支接合；
在 (7.1.41.3) 或 (7.1.41.6) 中，
删去 $D$，并把 $D$ 上的标记点
移至原来与它相连的分支。
重复此过程，直至不再有不稳定分支，
便得到 $\operatorname{Forget}_{\rho,\rho'}$。

\par\medskip\noindent
\begin{minipage}[t]{0.32\textwidth}\centering
\setlength{\unitlength}{0.67mm}
\begin{picture}(50,61)
\put(25,46){\circle{13}}
\put(25,30){\circle{20}}
\put(25,39){\circle*{1}}
\put(34,49){$\scriptstyle D$}
\put(25,18){\vector(0,-1){7}}
\put(25,6){\circle{18}}
\end{picture}
\par{\renewcommand{\thefigure}{7.1.8}\refstepcounter{figure}
\normalfont\bfseries 图~\thefigure\label{fig:7-1-8}}
\par{\small\textup{(7.1.41.1)}}
\end{minipage}\hfill
\begin{minipage}[t]{0.32\textwidth}\centering
\setlength{\unitlength}{0.67mm}
\begin{picture}(50,61)
\put(12,53){\circle{9}}
\put(9,53){$\scriptstyle +$}
\put(12,41){\circle{12}}
\put(12,27){\circle{17}}
\put(12,47){\circle*{1}}
\put(12,35){\circle*{1}}
\put(19,41){$\scriptstyle D$}
\put(23,34){\vector(1,0){5}}
\put(39,42){\circle{9}}
\put(36,42){$\scriptstyle +$}
\put(39,27){\circle{17}}
\put(39,35){\circle*{1}}
\end{picture}
\par{\renewcommand{\thefigure}{7.1.9}\refstepcounter{figure}
\normalfont\bfseries 图~\thefigure\label{fig:7-1-9}}
\par{\small\textup{(7.1.41.2)}}
\end{minipage}\hfill
\begin{minipage}[t]{0.32\textwidth}\centering
\setlength{\unitlength}{0.67mm}
\begin{picture}(50,61)
\put(25,46){\circle{13}}
\put(25,30){\circle{20}}
\put(25,39){\circle*{1}}
\put(25,46){\circle*{1}}
\put(22,50){$\scriptstyle z^+$}
\put(34,47){$\scriptstyle D$}
\put(18,29){$\scriptstyle\times$}
\put(29,29){$\scriptstyle\times$}
\put(25,18){\vector(0,-1){7}}
\put(25,6){\circle{18}}
\put(25,7){$\scriptstyle z^+$}
\put(18,5){$\scriptstyle\times$}
\put(29,5){$\scriptstyle\times$}
\end{picture}
\par{\renewcommand{\thefigure}{7.1.10}\refstepcounter{figure}
\normalfont\bfseries 图~\thefigure\label{fig:7-1-10}}
\par{\small\textup{(7.1.41.3)}}
\end{minipage}\par\medskip

\noindent
\begin{minipage}[t]{0.32\textwidth}\centering
\setlength{\unitlength}{0.67mm}
\begin{picture}(50,61)
\put(25,45){\circle{14}}
\put(25,29){\circle{21}}
\put(25,38){\circle*{1}}
\put(35,47){$\scriptstyle D$}
\put(35,29){$\scriptstyle z_0$}
\put(8,29){$\scriptstyle z_1$}
\put(12,18){$\scriptstyle z_2$}
\put(31,18){$\scriptstyle z_3$}
\put(25,17){\vector(0,-1){6}}
\put(25,6){\circle{18}}
\put(34,6){$\scriptstyle z_0$}
\put(8,6){$\scriptstyle z_1$}
\put(13,0){$\scriptstyle z_2$}
\put(31,0){$\scriptstyle z_3$}
\end{picture}
\par{\renewcommand{\thefigure}{7.1.11}\refstepcounter{figure}
\normalfont\bfseries 图~\thefigure\label{fig:7-1-11}}
\par{\small\textup{(7.1.41.4)}}
\end{minipage}\hfill
\begin{minipage}[t]{0.32\textwidth}\centering
\setlength{\unitlength}{0.67mm}
\begin{picture}(50,61)
\put(12,53){\circle{9}}
\put(12,41){\circle{12}}
\put(12,27){\circle{17}}
\put(12,47){\circle*{1}}
\put(12,35){\circle*{1}}
\put(19,41){$\scriptstyle D$}
\put(3,55){$\scriptstyle z_2$}
\put(16,55){$\scriptstyle z_1$}
\put(2,24){$\scriptstyle z_3$}
\put(16,24){$\scriptstyle z_0$}
\put(23,34){\vector(1,0){5}}
\put(39,42){\circle{9}}
\put(39,27){\circle{17}}
\put(39,35){\circle*{1}}
\put(30,45){$\scriptstyle z_2$}
\put(42,45){$\scriptstyle z_1$}
\put(30,23){$\scriptstyle z_3$}
\put(43,23){$\scriptstyle z_0$}
\end{picture}
\par{\renewcommand{\thefigure}{7.1.12}\refstepcounter{figure}
\normalfont\bfseries 图~\thefigure\label{fig:7-1-12}}
\par{\small\textup{(7.1.41.5)}}
\end{minipage}\hfill
\begin{minipage}[t]{0.32\textwidth}\centering
\setlength{\unitlength}{0.67mm}
\begin{picture}(50,61)
\put(25,47){\circle{20}}
\put(25,30){\circle{14}}
\put(25,37){\circle*{1}}
\put(17,56){$\scriptstyle z_2$}
\put(30,56){$\scriptstyle z_1$}
\put(36,47){$\scriptstyle z_0$}
\put(35,30){$\scriptstyle D$}
\put(16,23){$\scriptstyle z_3$}
\put(25,21){\vector(0,-1){5}}
\put(25,7){\circle{18}}
\put(17,15){$\scriptstyle z_2$}
\put(30,15){$\scriptstyle z_1$}
\put(35,7){$\scriptstyle z_0$}
\put(16,1){$\scriptstyle z_3$}
\end{picture}
\par{\renewcommand{\thefigure}{7.1.13}\refstepcounter{figure}
\normalfont\bfseries 图~\thefigure\label{fig:7-1-13}}
\par{\small\textup{(7.1.41.6)}}
\end{minipage}\par\medskip

\begin{definition}{7.1.42}\label{def:7-1-42}
设 $k\ge0$，且
\[
\bigl((\Sigma_i,\vec z_i,\vec z_i^{\, +}),w_i\bigr),\quad
\bigl((\Sigma_\infty,\vec z_\infty,\vec z_\infty^{\, +}),w_\infty\bigr)
\in\mathcal M^{\mathrm{main}}_{k+1,m}(L;\beta).
\]
称
\[
\lim_{i\to\infty}
\bigl((\Sigma_i,\vec z_i,\vec z_i^{\, +}),w_i\bigr)
=
\bigl((\Sigma_\infty,\vec z_\infty,\vec z_\infty^{\, +}),w_\infty\bigr),
\]
如果存在 $k'\ge k$、$m'\ge m$
以及保序单射
$\rho\colon\{1,\ldots,k\}\to\{1,\ldots,k'\}$
与单射
$\rho'\colon\{1,\ldots,m\}\to\{1,\ldots,m'\}$
\translatornote{原书在定义 7.1.42 仍称
$k',m',\rho,\rho'$“如上”，
但上文的箭头是从较多标记点的
$\mathcal M_{k+1,m}$ 遗忘到较少标记点的
$\mathcal M_{k'+1,m'}$；
本定义要把稳定化后的
$\mathcal M_{k'+1,m'}$ 遗忘回
$\mathcal M_{k+1,m}$，
故此处必须反向取单射及不等式。
见疑误 A-0426。}
及
\[
\bigl((\Sigma_i,\vec z_i,\vec z_i^{\, +})^\sim,\widetilde w_i\bigr),
\quad
\bigl((\Sigma_\infty,\vec z_\infty,\vec z_\infty^{\, +})^\sim,
\widetilde w_\infty\bigr)
\in\mathcal M^{\mathrm{main}}_{k'+1,m'}(L;\beta),
\]
使得
\[
\lim_{\mathrm{s},\,i\to\infty}
\bigl((\Sigma_i,\vec z_i,\vec z_i^{\, +})^\sim,\widetilde w_i\bigr)
=
\bigl((\Sigma_\infty,\vec z_\infty,\vec z_\infty^{\, +})^\sim,
\widetilde w_\infty\bigr),
\]
并且
\[
\begin{aligned}
\operatorname{Forget}_{\rho,\rho'}
\bigl((\Sigma_i,\vec z_i,\vec z_i^{\, +})^\sim,\widetilde w_i\bigr)
&=\bigl((\Sigma_i,\vec z_i,\vec z_i^{\, +}),w_i\bigr),\\
\operatorname{Forget}_{\rho,\rho'}
\bigl((\Sigma_\infty,\vec z_\infty,\vec z_\infty^{\, +})^\sim,
\widetilde w_\infty\bigr)
&=\bigl((\Sigma_\infty,\vec z_\infty,\vec z_\infty^{\, +}),w_\infty\bigr).
\end{aligned}
\]
\end{definition}

注意，为定义
$\mathcal M^{\mathrm{main}}_{k+1}(L;\beta)
=\mathcal M^{\mathrm{main}}_{k+1,0}(L;\beta)$
上的拓扑，我们使用了 $m>0$ 时的
$\mathcal M^{\mathrm{main}}_{k+1,m}(L;\beta)$。
由此将
$\mathcal M^{\mathrm{main}}_{k+1,m}(L;\beta)$
赋予可度量化的空间结构。
原书给出的理由是：该空间有有限分层，
每一层都是某个巴拿赫流形的子集\translatornote{仅有“有限分层”和“每层为巴拿赫流形的子集”
并不足以单独推出整体拓扑可度量；
仍须证明层间的拓扑相容性。
这里忠实保留原书论证，见疑误 A-0425。}。
上述事实也用来说明，
可以通过指定序列的极限而不是开集来定义拓扑。

\begin{theorem}{7.1.43}\label{thm:7-1-43}
$\mathcal M^{\mathrm{main}}_{k+1,m}(L;\beta)$
是紧豪斯多夫空间\liftIndex{I231}。
\end{theorem}

\begin{proof}
紧性在文献中已得到充分确立，
特别见
\textbf{[Grom85, Pan94, Ye94, IzSh00]}。
更具体地说，由于这里对拓扑的定义
与 \textbf{[FuOn99II]} 采取相同路线，
可利用这些文献中的分析结果，
仿照 \textbf{[FuOn99II]} 第~10、11 节
证明紧性。
伪全纯曲线模空间的紧性已是成熟结果，
故略去细节；
不过，文献中并不容易找到
伪全纯圆盘模空间拓扑的正确定义。
豪斯多夫性的证明与
\textbf{[FuOn99II]} 引理~10.4 相同。
\end{proof}
